\chapter[Where Elements Come From and Go]{Where Elements\protect\\ Come From and Go}

\begin{kaichang}
Imagine two things.

The first is a gold ring---perhaps on an older relative's finger. Gold does not rust or rot; leave it for centuries and it still gleams yellow. But where did Earth's gold come from? Once the mines are exhausted, can we ``grow another crop,'' as we grow grain? Can we ``synthesize'' more, as factories produce plastic?

The second is a sheet of paper you have burned (imagine it only; do not try it). Flames sweep over it, leaving ash and a wisp of smoke. Where did the paper go? Some people say it ``burned away.'' Did it really disappear?

``Where did it come from?'' and ``Where did it go?'' These sound like ordinary questions, but answering them takes us across the grandest scales in this book: from the first three minutes after the universe began, to the hearts of stars, to the soil beneath your feet and the rice in your bowl. By the end of this chapter, you will know that your gold ring's origins are far more romantic than you imagined.
\end{kaichang}

\section{Taking Stock: The Elements Around Us}

Before asking ``where did they come from,'' let us ask ``where are they now?'' We can take an ``element inventory'' of several familiar places, ranking the leaders by percentage mass:

\begin{table}[htbp]
\centering
\begin{tabular}{lllll}
\toprule
Rank & Crust & Seawater, incl. water & Atmosphere & Human body \\
\midrule
1st & O (about 46\%) & O (about 86\%) & N (about 78\%) & O (about 65\%) \\
2nd & Si (about 28\%) & H (about 11\%) & O (about 21\%) & C (about 18\%) \\
3rd & Al (about 8\%) & Cl, Na, etc. & Ar (about 0.9\%) & H (about 10\%) \\
4th & Fe (about 5\%) & Dozens more & Carbon dioxide, etc. & N (about 3\%) \\
\bottomrule
\end{tabular}
\end{table}

There are several surprises here. The most abundant element in Earth's crust is neither iron nor carbon, but oxygen: the main ingredients of sand, rocks, and clay have frameworks built from silicon and oxygen. Oxygen also takes first place by mass in the human body (because we are about 60 to 70 percent water), while carbon, the ``framework element of life,'' comes only second. The atmosphere is nitrogen's territory; oxygen accounts for only a fifth.

Remember, too, that these are averages. Your phone contains dozens of elements, including gallium, indium, and tantalum, which make up only a few parts per million of the crust. Bacteria near deep-sea hydrothermal vents contain far more molybdenum and tungsten than land organisms. Elements are distributed extremely unevenly: in some places they accumulate into ore deposits; elsewhere they are so dilute that a tonne of ore yields only a few grams of metal. This ``unevenness'' will become a major issue later.

\section{Chemical Reactions Cannot Change Atomic Nuclei}

Chapter 7 explained that an element's identity is its number of nuclear protons: 8 means oxygen, 79 means gold, with no room for one more or one fewer. Chapters 9 through 11 explained that every ``plot twist'' in a chemical reaction---forming bonds, breaking bonds, gaining or losing electrons---involves electrons outside the nucleus. Throughout, the nucleus stays backstage, unchanged.

Put those facts together, and we get this chapter's first rule:

\begin{qiaoliang}
\textbf{In every ordinary chemical reaction, the types of elements and the total number of atoms remain unchanged.} However many carbon atoms there were before the reaction, there are just as many afterward; they have only changed neighbors and arrangements. To turn one element into another, you must rewrite the nucleus itself. That is a nuclear reaction, requiring hundreds of thousands to millions of times more energy than chemical reactions---nothing a fire or a bottle of chemicals can accomplish.
\end{qiaoliang}

This gives us half the answer to ``where did the burned paper go?'' Not one of its carbon, hydrogen, or oxygen atoms disappeared. Most joined up with oxygen gas, becoming carbon dioxide and water vapor that dispersed into the air; a small fraction remained in the ash. It seems to have ``gone'' because the products changed from a solid you could see and touch into invisible gases. The visible form disappeared, not the atoms.

Everyday life offers ways to test this rule, as long as you include the ``hidden parts'' in your accounts. A rusty nail becomes heavier, but the extra mass did not appear from nowhere: it is oxygen ``captured'' from the air. Silver jewelry darkens because silver joins trace hydrogen sulfide in the air to form silver sulfide; no silver atoms are lost---they simply change partners. Some carbon in the rice you eat becomes the carbon dioxide you exhale, while some is built into your bones and muscles. \textbf{Not a single atom is ``manufactured'' inside you}; you are merely a busy transfer station for atoms. This also explains why we must obtain ``minerals'' such as iron, zinc, and iodine from food: the body can assemble molecules, but cannot conjure up elements. Missing elements must be supplied as ready-made atoms from outside.

The following experiment lets you ``catch'' these atoms that refuse to disappear.

\begin{shiyan}
\textbf{Conservation in a Sealed Bag}

Materials: a sealable transparent plastic bag (a zip-top food bag is best), a small spoonful of baking soda (sodium bicarbonate), a small cup of white vinegar, a paper towel, and a digital kitchen scale (a resolution of 1 gram is enough).

Procedure:
\begin{enumerate}[itemsep=1pt]
\item Pour the baking soda onto the paper towel, wrap it into a small packet, and place it in the plastic bag.
\item Pour the vinegar into the bag, taking care not to let it touch the packet yet.
\item Squeeze out most of the air, seal the bag tightly, place the whole bag on the scale, and record its mass.
\item Without moving the scale, pinch the packet open through the bag and jiggle it to let the baking soda react with the vinegar. Watch the bag slowly inflate, with bubbles filling it.
\item Once the reaction ends and the bag has inflated, read the mass again.
\end{enumerate}

What to observe: the scale's reading is almost unchanged (the difference is usually smaller than the scale's own uncertainty). A vigorous reaction clearly occurred and produced plenty of gas, yet the total mass stayed the same, because all the carbon dioxide molecules remained trapped inside. If you open the bag and weigh it again, the mass decreases---not because atoms disappeared, but because the gas ``escaped from prison.''

Safety note: this experiment is very safe, but work beside a sink in case the bag leaks and sprays vinegar solution. Do not replace the plastic bag with a sealed glass bottle; expanding gas could shatter the glass.
\end{shiyan}

\section{Who Does Rewrite the Elements in Atomic Nuclei?}

Chemical reactions cannot create gold or destroy carbon. So our original question returns unchanged: where did the world's hundred-plus elements come from? The answer is that they were born in three batches, in three completely different ``factories.''

\textbf{Factory one: the Big Bang.} Chapter 8 mentioned that the universe began about 13.8 billion years ago. The newborn universe was unimaginably hot; only as it expanded and cooled could protons and neutrons ``freeze out.'' During roughly the first three minutes, a small fraction happened to collide and stick together as helium nuclei, with a tiny amount forming lithium. Then the universe became too cool and diffuse for further fusion. The first batch of elements was therefore exceedingly simple: about three-quarters hydrogen and one-quarter helium by mass, plus a trace of lithium. Almost every hydrogen atom in your body was born in those three minutes and is 13.8 billion years old.

\textbf{Factory two: stars.} The Big Bang could not make carbon, oxygen, or iron---without which there would be no planets or life. These elements had to wait hundreds of millions of years, until gravity compressed clouds of hydrogen and helium into stars. Stars shine through nuclear fusion: at temperatures of tens of millions of degrees and immense pressures in their cores, hydrogen nuclei fuse into helium. The released energy resists gravity's compression, allowing steady burning for billions of years. After the hydrogen is exhausted, the core contracts and heats up further, and helium begins fusing into carbon and oxygen. Larger stars can keep burning successive layers: carbon into neon and magnesium, then into silicon, and finally into iron. Like an onion, an aging star has a different ``fuel'' in each internal layer.

This assembly line has a counterintuitive schedule: the later the ``fuel,'' the faster it burns. The Sun has enough hydrogen for about 10 billion years. Even massive stars burn hydrogen for millions to tens of millions of years, but their helium stage lasts only hundreds of thousands of years, their carbon stage a few hundred years, and their silicon stage---only about a day. It is as though a factory takes its time over the first few operations, rushes through the final one in a day, and then shuts down with a crash. Hence the universe is piled high with light elements, while elements above iron remain a tiny fraction: hydrogen and helium account for more than 98\% of visible matter, and all the other hundred-plus elements together account for less than 2\%.

But the road ends at iron. We glimpsed the reason in Chapter 8; here is the fuller explanation:

\begin{qiaoliang}
Inside a nucleus, protons repel one another, and the strong nuclear force holds them together. Bringing separate nucleons together into a nucleus releases energy, called ``binding energy.'' Per nucleon, binding energy peaks near iron: an iron nucleus is in the ``tightest, most comfortable'' state. Thus fusion of elements lighter than iron lowers total energy and releases energy, whereas fusion of elements heavier than iron requires an energy investment. Iron is the final stop of stellar nuclear burning: once a star's core has burned into iron, no more energy can be squeezed out to oppose gravity, and a massive star collapses and explodes.
\end{qiaoliang}

\textbf{Factory three: stellar deaths and collisions.} Elements heavier than iron---silver, gold, platinum, uranium, and iodine---can only be made in more violent settings. One is a supernova explosion at the death of a massive star: the enormous energy released during collapse smashes iron nuclei apart and reshapes them, forcing nucleons into heavier nuclei. Another is a pair of neutron stars (dense stellar remnants, a teaspoonful of which weighs hundreds of millions of tonnes) orbiting each other and finally colliding. The collision showers surrounding nuclei with neutrons; a nucleus swallows dozens at a time and, after a series of decays, settles into a heavy element such as gold or platinum. Every gold atom in your ring came from a cosmic pileup of this kind, before the solar system was born.

\section{Writing the Three Factories in Symbols}

Now it is time for symbols. Nuclear and chemical reaction equations look similar, but follow different rules: chemical equations conserve the types of atoms on either side; nuclear equations conserve the total number of protons and the sum of mass numbers, while changing the elements is precisely the point. Consider these three:

Hydrogen fusion (the Sun's main energy source, in simplified form):
\[4\,\chem{^{1}_{1}H} \rightarrow \chem{^{4}_{2}He} + \text{energy}\]

The birth of carbon (the ``triple-alpha process'' in red giants):
\[3\,\chem{^{4}_{2}He} \rightarrow \chem{^{12}_{6}C} + \text{energy}\]

One route to gold (a nucleus decays after capturing neutrons in a neutron-star merger):
\[\chem{^{197}_{80}Hg} \rightarrow \chem{^{197}_{79}Au} + \chem{^{0}_{+1}e}\]

Notice the third equation: mercury (number 80) really becomes gold (number 79) through decay. Medieval alchemists' dream is not physically impossible. So why not make gold and get rich? The misconception box in this chapter will explain.

Let us also calculate what it means to say ``every atom has an extraordinary history.'' How many gold atoms are in a \ud{5}{g} ring? Gold's molar mass is about \ud{197}{g\cdot mol^{-1}}, so $5/197 \approx 0.025\ \mathrm{mol}$. Multiplying by Avogadro's constant:
\[0.025 \times 6.02\times 10^{23} \approx 1.5\times 10^{22}\ \text{gold atoms}\]
One hundred and fifty quintillion atoms, each forged in the blazing collision of neutron stars billions of years ago, then carried through gas clouds, meteorites, and magma for billions more years before coming to rest in this ring. Romantic? This is calculation, not poetic license.

\section{How Do We Know That Stars Make Elements?}

We have never entered a star's core. How can we be so sure elements are made there? Because we have two kinds of compelling evidence: one from light, and one from ``predictions made visible.''

\begin{zhengju}
\textbf{Clue one: an element discovered on the Sun before it was found on Earth.} During a total solar eclipse in 1868, astronomers pointed spectroscopes at prominences along the Sun's edge and saw a bright yellow spectral line matching none of the elements then known on Earth. Someone made a bold announcement: the Sun contained an element not yet discovered on Earth, named helium after the Greek word for ``sun.'' This was a risky claim: what if the line came from an instrumental error or an unknown state of a known element? Debate continued for nearly thirty years. Then, in 1895, the chemist Ramsay found the same yellow line in gas released by heating a terrestrial uranium mineral. Helium really existed, and Earth had it too. This established spectroscopy's standing: \textbf{each element's emission lines are unique, like a barcode}, and light is a composition report from distant matter. Astronomers could now read that the Sun was about three-quarters hydrogen and one-quarter helium, matching the Big Bang's ``first-batch recipe.'' Carbon, oxygen, sodium, and iron lines in aging stars' spectra strengthen with age, matching the story of ``stars making elements as they burn.''

\textbf{Clue two: a prediction fulfilled before the world's eyes.} In 1957, the Burbidges, Fowler, and Hoyle published a paper systematically setting out the theory of heavy-element synthesis in stars and supernovae. It predicted that the heaviest elements, such as gold and platinum, required rapid neutron capture in a ``neutron downpour,'' and that colliding neutron stars should be among the ideal settings. Sixty years of waiting followed. In August 2017, gravitational-wave detectors caught the ripples from two neutron stars colliding about 130 million light-years away. Almost simultaneously, dozens of telescopes worldwide turned toward the same direction, seeing a rapidly expanding, fading afterglow in visible and infrared light: a ``kilonova.'' Its spectral evolution matched calculations for ``large quantities of heavy elements forming and decaying.'' Estimates suggest that this single collision ejected several hundred Earth masses of gold. From spectral barcodes to gravitational waves, two entirely independent chains of evidence lead to the same conclusion: stars and their deaths are the universe's element forges.
\end{zhengju}

\section{Elements Stay on Earth, but Keep Moving House}

Elements do not appear or disappear from nowhere, but neither do they sit still. They continually move among rocks, oceans, the atmosphere, and living things. Scientists call these routes ``cycles.'' Let us look at three especially important to life:

\begin{itemize}[itemsep=2pt]
\item \textbf{The carbon cycle}: plants fix carbon dioxide into sugars and lignin through photosynthesis; animals eat plants; and respiration, death, and decay return carbon to the atmosphere and soil. Some carbon settles on the seafloor and becomes limestone, waiting for crustal movements to lift it back to the surface. A carbon atom's ``commute'' can last days (through respiration) or over a hundred million years (through limestone). Over the past two centuries, humans have rapidly burned carbon locked in coal and oil for hundreds of millions of years back into the atmosphere, like moving a whole load of freight from the slow lane to the fast lane at once. The carbon atoms are unchanged, but their atmospheric concentration has changed, and the climate with it. Chapter 54, on photosynthesis, returns to this cycle's ``entrance.''
\item \textbf{The nitrogen cycle}: four-fifths of the air is nitrogen gas, but its molecules are so stable that most organisms cannot use them. Only a few bacteria and lightning can ``fix'' nitrogen into usable nitrogen compounds. Plants absorb these, animals eat the plants, and microorganisms decompose their remains and wastes, returning some nitrogen to the sky as nitrogen gas. Over a century ago, humans invented industrial nitrogen fixation (the Haber process), using high temperatures and pressures to combine nitrogen and hydrogen into ammonia. Today, the crops feeding about half the world's population depend on nitrogen-fixation factories. This is one of the most powerful examples of chemistry changing civilization, discussed in detail in Chapter 14.
\item \textbf{The phosphorus cycle}: phosphorus has no gaseous form and moves slowly among rocks, soil, water, and organisms. Weathering releases phosphates from rocks; plants absorb them; they pass along food chains; and finally they settle underwater in dead organisms, waiting tens of millions of years for geological processes to bring them back. Phosphorus travels on something close to a ``one-way ticket,'' so farmland needs phosphate fertilizer. Meanwhile, domestic wastewater carries large amounts of phosphorus into lakes, promoting runaway algal growth and oxygen depletion (eutrophication).
\end{itemize}

Notice the common pattern: \textbf{none of these cycles creates or destroys elements}. A ``cycle'' is simply a collection of age-old atoms changing partners and addresses. Almost all the energy driving their moves ultimately comes from the Sun. (The metabolism chapters from Chapter 54 onward explain how life uses this energy to organize atoms' movements into the business of ``being alive.'')

\section{Why Shortages If Elements Never Disappear?}

You may want to object: if elements never disappear, why do the news reports keep warning of ``shortages'' of lithium, cobalt, and rare earths? Good question. This is exactly where the chapter's ideas meet industry.

A shortage never means ``the atoms are gone.'' It means two things. First, \textbf{concentration}: a phone's lithium battery contains only a few grams of lithium, and an electric car a few kilograms. After use, that lithium has not disappeared; it is scattered among hundreds of millions of discarded batteries, mixed with plastics, aluminum, and electrolytes. Gathering and purifying scattered elements costs energy and money. Dilution is easy; concentration is hard---the second law of thermodynamics, coming in Chapter 27, shows itself here in resource management. Second, \textbf{distribution}: rare earths are not particularly rare in the crust, but deposits concentrated enough to mine are scarce. More than half of cobalt reserves are concentrated in the Democratic Republic of the Congo, while highly concentrated lithium brines occur mainly in South America's ``Lithium Triangle'' and China's Qinghai--Tibet Plateau. Elements cycle globally, but ore deposits are extremely selective about location, turning supply chains into a geopolitical issue.

The solution follows from this chapter's logic: since atoms persist, \textbf{recycling is ``urban mining.''} A tonne of discarded phone circuit boards can contain dozens of times more gold than a tonne of gold ore. Japan once collected about 80,000 tonnes of discarded small electronics for Tokyo Olympic medals, extracting about \ud{32}{kg} of gold, \ud{3500}{kg} of silver, and \ud{2200}{kg} of copper. Keeping elements circulating in the economy for a few more rounds often takes less energy and causes less pollution than chasing new deposits deep underground.

\section{Where This Rule Breaks Down}

Let us be honest about the boundaries. ``Elements are neither created nor destroyed'' holds only at the scale of \textbf{chemical and biological processes}. Beyond that, qualifications are needed:

\begin{itemize}[itemsep=2pt]
\item \textbf{Nuclear reactions do not follow this rule}. The radioactive decay in Chapter 8 slowly turns uranium into lead; hospital nuclear medicine transforms elements daily; and the Sun is turning hydrogen into helium right now. These processes operate at energy scales several orders of magnitude removed from ordinary life and chemistry.
\item \textbf{The ``inventory'' is not eternal}. The elemental proportions of the crust and atmosphere change over geological time. More than two billion years ago, the atmosphere held almost no oxygen. Cyanobacteria's photosynthesis took hundreds of millions of years to ``make'' it (the oxygen atoms were not new; their release from water was). Radioactive elements such as uranium-238, meanwhile, decrease irreversibly. With a half-life of about 4.5 billion years, nearly half the solar system's uranium has already been used up.
\item \textbf{An extremely slight ``mass--element'' conversion}. Any energy-releasing reaction leaves products with slightly less total mass (see the challenge box). In chemical reactions, however, the difference is too tiny for a balance ever to detect, so ``atom conservation'' is an excellent approximation in chemistry. Calling it an ``approximation'' does not mean it is wrong; it means we know the range of validity printed in its instruction manual.
\end{itemize}

\begin{wuqu}
\textbf{``If nuclear reactions can change elements, why not make gold and get rich?''}

The idea works; the economics do not. Laboratories really have used particle accelerators to turn bismuth or mercury into gold. In 1980, a team at America's Lawrence Berkeley Laboratory bombarded a bismuth target with an accelerator and obtained tiny quantities of gold. The price? Counting electricity and equipment depreciation, making 1 gram costs millions of times the gold's value. The products also often contain radioactive isotopes that must decay to safe levels before handling. By contrast, historical alchemists spent centuries with furnaces and chemical concoctions without making a single gold atom: their methods were all chemical, and chemistry cannot touch atomic nuclei. This misconception is worth remembering because it sharply separates two things: \textbf{between ``physically possible'' and ``economically viable'' lies an enormous gulf of energy, cost, and scale}. Nuclear physics has breached one corner of element conservation, but economics has built ten walls in its place.
\end{wuqu}

\section{Back to the Gold Ring and the Paper}

Now we can fulfill our two opening promises.

The gold ring: each atom was born in a neutron-star collision (or supernova) before the solar system formed, joining a gas cloud in the ejected debris. About 4.6 billion years ago, that cloud contracted into the solar system, and the gold accompanied other heavy elements into the forming Earth. While the early Earth was molten, most gold sank toward the core. A small fraction remained in the crust and was later concentrated into ore by hydrothermal fluids. Miners dug it up, and craftspeople shaped it into a ring. Gold cannot be ``grown,'' nor can chemical reactions add or remove any of it. Earth's gold is essentially that original batch: each use disperses some, and recycling is the only way to gather it back together.

The burned paper: most of its carbon and hydrogen became \chem{CO_2} and \chem{H_2O} and entered the air. You can continue the story yourself. A carbon dioxide molecule might be captured by a leaf next week and become glucose. It might dissolve in seawater, be made into a plankton organism's calcium carbonate shell, settle to the seafloor, become limestone millions of years later, then be blasted apart, fired into cement, and built into a wall. Atoms do not die; they are always hurrying to their next engagement.

\section{Looking Ahead: Same Atoms, New Partners}

By focusing on nuclei, this chapter has given us an almost static picture: elements are conserved, and atoms endure. Yet everyday substances change constantly---paper burns, iron rusts, and food becomes part of you. The entire secret lies not in atoms' births and deaths, but in rearranging their \textbf{connections}: who partners with whom, how strongly they pull, and how the energy accounts balance during rearrangement. Starting next chapter, our camera moves from atomic nuclei to the spaces between atoms. In Chapter 17, we will take apart one of chemistry's most misunderstood concepts: the ``chemical bond.''

\begin{tiaozhan}
\textbf{Where Does a Star Get So Much Energy? A Mass Balance.} Fusion energy comes from mass itself: when 4 protons fuse into 1 helium nucleus, the products have about $0.7\%$ less mass than the reactants. Using $E=mc^2$, fusing 1 gram of hydrogen loses about \ud{0.007}{g} of mass and releases about $0.007\times 10^{-3}\ \mathrm{kg}\times (3\times 10^{8}\ \mathrm{m/s})^2 \approx 6\times 10^{11}\ \mathrm{J}$---equivalent to burning about 20 tonnes of coal. Each second, the Sun burns about 600 million tonnes of hydrogen into helium and ``loses'' over 4 million tonnes of mass. The astronomical size of the speed of light squared is what lets it keep burning after 4.6 billion years (using only a few percent of its total hydrogen so far). That enormous $c^2$ also explains why mass changes in chemical reactions (on the order of a few parts per billion) are completely unmeasurable by weighing. You can skip this box without losing the main thread, but it connects Chapter 8's binding energy, this chapter's ``iron is the endpoint,'' and the energy accounts in the next part.
\end{tiaozhan}

\section*{Try It}

\begin{sikaoti}
\item Draw a flowchart of the ``three element factories'': use three boxes (Big Bang, stellar fusion, stellar death/collision) and arrows to show which factory produces hydrogen, helium, carbon, oxygen, iron, and gold. Add a note that box sizes do not represent actual scales of time or space.
\item After a piece of firewood burns, its ash weighs much less than the wood. A classmate says, ``This shows conservation of mass does not apply to burning.'' Refute the claim in particle-level terms, and design an improved test using only kitchen materials to support your argument.
\item Look up the approximate \ud{500}{mL} of air exchanged in one breath. Using the fact that the atmosphere is about 78\% nitrogen, estimate how many moles of nitrogen you inhale each day (1 mole of gas occupies about \ud{24}{L} at room temperature and ordinary pressure). Most of this nitrogen is ``exhaled unchanged.'' What does that reveal about the nitrogen cycle?
\item Following the gold-ring calculation, estimate how many iron atoms are in the roughly \ud{4}{g} of iron in your body (mostly in hemoglobin). Explain which ``factory'' most likely produced them. (Hint: iron's special position in element synthesis.)
\item A supplement advertisement claims to contain ``special active elements that the body cannot synthesize and must receive as extra supplements.'' Use this chapter to identify the reasonable part and the conceptual sleight of hand. (Is it the elements themselves that the body cannot synthesize, or particular molecules containing them?)
\item Draw a route for ``a carbon atom's ten-thousand-year journey'': start in the atmosphere, visit at least three stops among plants, animals, soil, and oceans, and return to the atmosphere. At each stop, label the atom's new ``partners'' (\chem{CO_2}, glucose, carbonate, and so on).
\end{sikaoti}
